\begin{figure}[h]
    \centering
    \vspace{-7mm}
    \begin{subfigure}{0.45\textwidth}
        \centering
        \includegraphics[width=\linewidth]{figs/all.pdf}
        \caption{Typical LLMs’ AgentBench performance (relative) against the best in each environment.}
    \end{subfigure}
    \hspace{3mm}
    \begin{subfigure}{0.45\textwidth}
        \centering
        \includegraphics[width=\linewidth]{figs/overal-bar.pdf}
        \caption{Overall scores of AgentBench across 8 environments. Dashed lines for two LLM types’ average.}
    \end{subfigure}
    \vspace{-3mm}
    \caption{An overview of LLMs on \model. While LLMs begin to manifest their proficiency in LLM-as-Agent, gaps between models and the distance toward practical usability are significant.}
    \label{fig:intro}
    \vspace{-6mm}
\end{figure}

\section{Introduction}

Intelligent agents and autonomous entities~\citep{Searle1970SpeechAA,Maes1994AgentsTR,wooldridge1995intelligent} that are capable of decision-making and action execution in particular environments have been key concepts of artificial intelligence (AI) historically.
Notwithstanding substantial advancements in deep learning algorithms applied in both computer vision and natural language processing (NLP), their potential for developing efficient and practically usable assisting agents remains largely unexplored.

The advent of Large Language Models (LLMs)~\citep{GPT3,chowdhery2022palm,touvron2023llama}, such as GPT-4~\citep{openai2023gpt}, has brought plenty of new opportunities to this realm.
Through extensive alignment training~\citep{ouyang2022training,wei2022finetuned,sanh2022multitask}, LLMs have not only mastered traditional NLP tasks but also showcased an impressive ability to comprehend human intent and execute instructions.
This has spurred the development of various LLM-based applications for autonomous goal completion (like AutoGPT~\citep{richards2023auto}, BabyAGI~\citep{nakajima2023babyagi}, AgentGPT~\citep{agentgpt}) as well as LLM agents situated in social and game contexts~\citep{Park2023GenerativeAI,Wang2023VoyagerAO,Zhu2023GhostIT}, sparking substantial public interest and discussions. 

Despite these advancements, the lack of a systematic and standard benchmark to evaluate LLM-as-Agent presents a critical challenge. 
Historically, text-based game environments~\citep{osborne2022survey,cote2019textworld,hausknecht2020interactive,urbanek2019learning} have been employed for language agent evaluation.
But they often suffer from the limitation of closed, discrete action spaces, as well as their primarily narrow focus on models' commonsense grounding.
More recently, attempts on embodied agents~\citep{reed2022generalist,huang2022language,ahn2022can,li2022pre} have employed complicated multi-modal simulators based on games~\citep{kuttler2020nethack,fan2022minedojo}, GUI~\citep{shi2017world,toyama2021androidenv}, and indoor scenes~\citep{shen2021igibson,srivastava2022behavior}.
However, these simulators, despite their complexity, do not accurately reflect the practical use cases of LLMs, and their multi-modal nature creates a hurdle for the urgent evaluation of existing text-only LLMs.
Finally, most benchmarks now for agents focus on single environments and thus fail to provide a comprehensive overview of LLMs across diverse application scenarios.

\begin{figure}[t]
    \centering
    \vspace{-2mm}
    \includegraphics[width=.88\linewidth]{figs/agentbench.pdf}
    \vspace{-2mm}
    \caption{\model is the first systematic benchmark to evaluate LLM-as-Agent on a wide array of real-world challenges and 8 distinct environments. In total, \num LLMs are examined in this edition.}
    \label{fig:framework}
    \vspace{-8mm}
\end{figure}

To address these challenges, we introduce \model, a multi-dimensional benchmark designed to evaluate LLM-as-Agent across a spectrum of different environments.
\model encompasses eight distinct environments (Cf. Figure~\ref{fig:demo}, five out of eight are created for the first time), which could be categorized into three types of groundings: 

\begin{itemize}[leftmargin=1.5em,itemsep=0pt,parsep=0.2em,topsep=0.0em,partopsep=0.0em]
\item \textbf{Code:} Operating System, Database, Knowledge Graph~\citep{kgtool}
\item \textbf{Game:} Digital Card Game, Lateral Thinking Puzzles, House-Holding~\citep{shridhar2020alfworld}
\item \textbf{Web:} Web Shopping~\citep{yao2022webshop}, Web Browsing~\citep{deng2023mind2web}
\end{itemize}

All datasets, whether newly created or adapted from existing ones, are meticulously designed and reformulated to simulate interactive environments where text-only LLMs can operate as autonomous agents.
\model thus systematically evaluate an LLM's core abilities, including following instructions~\citep{ouyang2022training}, coding~\citep{chen2021evaluating}, knowledge acquisition~\citep{joshi2017triviaqa,talmor2019commonsenseqa}, logical reasoning~\citep{srivastava2023beyond}, and commonsense grounding~\citep{shridhar2020alfred}.
It serves as an ideal testbed for both LLM and agent evaluation.

In addition, we develop a unified evaluation toolkit for LLMs to operate on diverse customized agent tasks, thus enabling a comprehensive benchmarking of the LLM-as-Agent ability of \num different LLMs on \model, including both API-based and OSS models.
Our results reveal that top-tier models like GPT-4 are capable of handling a wide array of real-world tasks, indicating the potential for developing a potent, continuously learning agent.
However, we also note a significant performance gap between these top-tier models and their OSS competitors.
Despite the recent success of OSS LLMs and their competitive scores on several benchmarks~\citep{li2023alpacaeval,chen2021evaluating,cobbe2021training}, their performance on the challenging \model tasks lags considerably.
This underscores the necessity for additional efforts to enhance the learning abilities of OSS LLMs.

We identify portions of agent task failures in different environments and LLMs, unveiling the insufficient abilities of long-term reasoning, decision-making, and instruction following in existing LLMs.
Comparisons between different LLMs manifest that a proper strategy of introducing code training can help improve LLM-as-Agent. Alignment training over high-quality data (e.g., data generated by \texttt{gpt-4}) could also help improve LLM agents.
In summary, our contributions are:

\begin{itemize}[leftmargin=1.5em,itemsep=0pt,parsep=0.2em,topsep=0.0em,partopsep=0.0em]
\item We introduce the concept of evaluating LLMs as agents and present \model, a comprehensive benchmark to standardize the evaluation. It defines eight distinct environments of 3 types based on real-world scenarios, offering a practical testbed for LLMs' wide array of capabilities.
\item We perform a thorough evaluation of \num different LLMs using \model, uncovering a significant performance gap between leading API-based commercial LLMs and many OSS models that are no larger than 70B. We also quantitatively analyze the reasons for failures in existing LLM agents and highlight directions for improvement, such as improving instruction following, higher-quality alignment data. Also, we show that code training could be a double-edged sword, which improves some agent tasks while harms other.
\item To facilitate the evaluation of LLM-as-Agent, we have introduced an integrated toolkit grounded in the Server-Client architecture, focusing on modular and scalable design principles. This enables easy customization of model assessments for any LLMs using the HTTP protocol. Complemented by its associated datasets and environments, this toolkit is now openly accessible to the broader research community.
\end{itemize}


